\documentclass[journal]{IEEEtran}
\usepackage{cite}
\usepackage{amsfonts}
\usepackage{amssymb}
\usepackage{amsmath}
\usepackage{amsthm}
\usepackage{url}
\usepackage{enumitem}
\usepackage{array}
\usepackage{booktabs}
\usepackage[utf8]{inputenc}
\usepackage[english]{babel}

\newtheorem{definition}{Definition}
\newtheorem{assumption}{Assumption}
\newtheorem{lemma}{Lemma}
\newtheorem{corollary}{Corollary}
\newtheorem{law}{Law}
\newtheorem{invariant}{Invariant}
\newtheorem{proposition}{Proposition}
\newtheorem{obligation}{Obligation}

\newcommand{\Nm}{\mathcal{N}}
\newcommand{\Vl}{\mathcal{V}}
\newcommand{\dom}{\operatorname{dom}}
\newcommand{\rec}{\operatorname{rec}}
\newcommand{\wit}{\operatorname{wit}}
\newcommand{\name}{\operatorname{name}}
\newcommand{\obs}{\operatorname{observe}}
\newcommand{\rmin}{\operatorname{rmin}}
\newcommand{\stored}{\operatorname{stored}}
\newcommand{\req}{\operatorname{req}}
\newcommand{\Ok}{\mathsf{Ok}}
\newcommand{\Err}{\mathsf{Err}}
\newcommand{\Filed}{\mathsf{Filed}}
\newcommand{\Undec}{\mathsf{Undecodable}}
\newcommand{\Missing}{\mathsf{Missing}}
\newcommand{\Exhausted}{\mathsf{BudgetExhausted}}
\newcommand{\Volatile}{\mathsf{Vol}}
\newcommand{\Cached}{\mathsf{Cch}}
\newcommand{\Durable}{\mathsf{Dur}}
\newcommand{\Rebuild}{\mathsf{Reb}}
\newcommand{\Retained}{\mathsf{Ret}}
\newcommand{\Auth}{\mathsf{Aut}}
\newcommand{\Retire}{\textsc{Retire}}
\newcommand{\Keep}{\textsc{Keep}}
\newcommand{\Requeue}{\textsc{Requeue}}
\newcommand{\Defer}{\textsc{Defer}}
\newcommand{\U}{\mathcal{U}}
\newcommand{\K}{\mathcal{K}}
\newcommand{\W}{\mathcal{W}}
\newcommand{\issue}[1]{\#{#1}}
\newcommand{\jsq}[1]{\texttt{#1}}

\author{Ryan~J.~Kung%
\thanks{Ryan J. Kung can be reached at ryankung@ieee.org.
Manuscript created October 23, 2026.  This document specifies the durability model of the
local stores of Rings before 1.0 (\texttt{rings-core}, modules \texttt{storage},
\texttt{message::replay} and \texttt{swarm::inbox}; \texttt{rings-node}, module
\texttt{onion::entry\_guard}).  It is the law for issues \issue{911}, \issue{912},
\issue{915}, \issue{923} and \issue{936} of \texttt{RingsNetwork/rings}.}}
\title{\huge Durable Records:\\
One Presence and Survival Law\\
for the Local Stores of a Rings Node}
\date{}

\IEEEaftertitletext{%
\vspace{0.5\baselineskip}%
\noindent\hfill
\parbox{0.90\textwidth}{%
{\centering\bfseries Abstract\par}%
\vspace{0.4\baselineskip}%
\noindent
A Rings node keeps security state (transaction replay windows), privacy state (onion entry
guards) and custody state (relay mail) in local key-value stores on two platform families:
a native file system and a browser's IndexedDB.  Five open defects of these stores reduce to
two predicates: \emph{presence}, whether a record that should exist exists, and
\emph{survival}, whether a completed write outlives a crash.  We give one model for both.
A store's authority is a point of the product of a crash chain
$\Volatile\sqsubset\Cached\sqsubset\Durable$ and a meaning chain
$\Rebuild\sqsubset\Retained\sqsubset\Auth$; each point fixes what an absent record means and
which obligations the platform must discharge.  We choose a per-record presence marker as
the witness that separates a lost authoritative record from a never-written one, prove it
sound and complete for partial loss, and prove that no witness in the same failure domain
detects total loss.  We state the crash obligation per platform, verified against the
Chromium, Gecko and WebKit sources (only \jsq{strict} durability flushes in all three),
forbid eviction for authoritative records, and give a custody law that makes relay mail
survive a restart between admission and consumption, with the impossibility result that
bounds what any receiver can do for mail that has fallen behind live traffic.
}\hfill\null\par
\vspace{0.5\baselineskip}}

\begin{document}
\maketitle
\emergencystretch=1.5em

\section{Problem}\label{sec:problem}

The defects, each named by its issue:
\begin{description}[leftmargin=1.2em,style=nextline]
\item[\issue{915}] An absent native replay record is indistinguishable from a stream never
seen, so its stream restarts from $\mathsf{First}$; on non-unix targets a completed rename
can roll back.
\item[\issue{923}] The browser replay store evicts its least recently used row at its row
cap, so a retained stream can restart from $\mathsf{First}$.
\item[\issue{912}] Browser replay writes use the default durability hint, so a completed
admission can roll back after an OS crash.
\item[\issue{911}] The entry-guard store is disposable: a crash can lose or tear the guard
set, re-selecting the first hop.
\item[\issue{936}] Relay mail admitted into the replay window and queued, but not taken
before a restart, is refused as a replay on redelivery and lost.
\end{description}
\issue{915}, \issue{923} are presence defects; \issue{912}, \issue{911} are survival defects;
\issue{936} is a presence defect of a fact that lives in no store at all (that the queued mail
was never consumed).  Section~\ref{sec:model} fixes the model, Section~\ref{sec:lattice} the
authority lattice, Sections~\ref{sec:witness}--\ref{sec:custody} the five laws, and
Section~\ref{sec:plan} the order of implementation.

\section{Model}\label{sec:model}

\begin{definition}[Store]\label{def:store}
A store is a finite partial map $\sigma:\Nm\rightharpoonup\Vl$ from record names to values.
Its operations are $\mathrm{put}(n,v)$, $\mathrm{rm}(n)$ and $\mathrm{get}(n)$, each
returning $\Ok$ or $\Err$; $\mathrm{scan}$ returns one classification per name
(Section~\ref{sec:witness}).  For \texttt{FileStorage} the name of key $k$ is
$\name(k)=\mathrm{hex}(\mathrm{SHA1}(k))$; for \texttt{IdbStorage} it is $k$.
\end{definition}

\begin{definition}[History, completion]\label{def:history}
A history is $h=o_1\cdots o_k\,o_{k+1}^{?}$: operations $o_1..o_k$ returned, $o_{k+1}$ may
be in flight.  Let $\sigma_i=o_i\circ\cdots\circ o_1(\sigma_0)$, with $\sigma_{k+1}=\sigma_k$
when nothing is in flight.  An operation is \emph{completed} iff it returned $\Ok$.
\end{definition}

\begin{definition}[Crash, recovery image]\label{def:crash}
A crash ends the process (\emph{process crash}) or the operating system
(\emph{system crash}: kernel panic, power loss).  The recovery image $R_P(h)$ is the set of
maps the next open of the store can read on platform $P$ after a crash at the end of $h$.
\end{definition}

\begin{definition}[Loss events]\label{def:loss}
A loss event is a change of the store's physical state that no operation of the store made:
\begin{itemize}[leftmargin=1.2em]
\item $\lambda_{\mathrm{part}}(S)$, \emph{partial loss}: the records of $S\subsetneq\dom\sigma$
vanish (media error repaired by \texttt{fsck}, an operator removing or renaming one file, a
completed rename rolled back by a platform that does not discharge
Obligation~\ref{ob:survival});
\item $\lambda_{\mathrm{tot}}$, \emph{total loss}: the whole failure domain vanishes (the
store's directory, an IndexedDB database or origin, a process heap).
\end{itemize}
\end{definition}

Loss events are what the platform or the operator does, not what an adversary forges: a
party that can write the store can also write any witness of it, so nothing below
authenticates a store, and nothing models a state the program and platform cannot reach.

\subsection{Obligations}

\begin{obligation}[Atomicity, A]\label{ob:atomicity}
$\forall\sigma'\in R_P(h).\ \forall n.\ \sigma'(n)\in\{\sigma_i(n): i\le k+1\}\cup\{\bot\}$:
a record is never torn.
\end{obligation}

\begin{obligation}[Survival, D]\label{ob:survival}
$\forall\sigma'\in R_P(h).\ \forall n.\ \sigma'(n)\in\{\sigma_k(n),\sigma_{k+1}(n)\}$:
every completed write and removal survives a process or system crash, and the one in
flight is either whole or absent.  D implies A.
\end{obligation}

\begin{obligation}[No eviction, N]\label{ob:noevict}
$\mathrm{put}(n,v)$ on $\sigma$ yields $\sigma[n\mapsto v]$ or returns
$\Err(\Exhausted)$ with $\sigma$ unchanged; no operation removes a name it was not given.
\end{obligation}

\begin{obligation}[Witnessed presence, W]\label{ob:witness}
After any history and any partial loss $\lambda_{\mathrm{part}}(S)$, $\mathrm{scan}$ reports
every $n\in S\cap\dom\sigma_k$ as $\Missing(n)$, and reports $\Missing(n)$ only for an $n$
whose record some $\mathrm{put}$ made durable.
\end{obligation}

\begin{obligation}[Fail closed, F]\label{ob:failclosed}
A record that does not decode is reported and kept; only its owner or an operator removes it.
\end{obligation}

\section{The authority lattice}\label{sec:lattice}

\begin{definition}[Crash chain]\label{def:crashchain}
$C=\Volatile\sqsubset\Cached\sqsubset\Durable$ with obligations
$O(\Volatile)=\emptyset$ ($R_P(h)=\{\emptyset\}$ at every restart), $O(\Cached)=\emptyset$
(any $\sigma_i(n)$, $\bot$ or a torn record), $O(\Durable)=\{A,D\}$.
\end{definition}

\begin{definition}[Meaning chain]\label{def:meaningchain}
$M=\Rebuild\sqsubset\Retained\sqsubset\Auth$, read as what $\sigma'(n)=\bot$ means to the
owner:
\begin{itemize}[leftmargin=1.2em]
\item $\Rebuild$: $\bot$ is the owner's default; any record may be dropped;
\item $\Retained$: $\bot$ is ``never written, or retired by the owner''; a loss is a
liveness failure;
\item $\Auth$: $\bot$ is ``never written''; a loss is a safety failure.
\end{itemize}
$O(\Rebuild)=\emptyset$, $O(\Retained)=\{N\}$, $O(\Auth)=\{N,W,F\}$.
\end{definition}

\begin{definition}[Authority]\label{def:authority}
An authority is $a=(c,m)\in C\times M$, ordered componentwise, with
$O(c,m)=O(c)\cup O(m)$.  $C\times M$ is a lattice (a product of chains, meet and join
componentwise) and $O$ is monotone: $a\sqsubseteq b\Rightarrow O(a)\subseteq O(b)$.
\end{definition}

\begin{law}[Admissibility]\label{law:admissible}
Let $\req(\Rebuild)=\Volatile$ and $\req(\Retained)=\req(\Auth)=\Durable$.  A store may be
opened at $(c,m)$ only if $c\sqsupseteq\req(m)$.
\end{law}
\begin{proof}[Justification]
Under $c\sqsubset\Durable$ a crash may lose a completed write, so $\bot$ after recovery can
be a loss; a meaning above $\Rebuild$ asserts that $\bot$ is not a loss.
\end{proof}

The existing \texttt{RecordAuthority} is the two-point chain
$\mathsf{Disposable}=(\Cached,\Rebuild)\sqsubset\mathsf{Authoritative}=(\Durable,\Auth)$
(\#909, \#910, \#914).  The law adds two points:
$\mathsf{DurableCache}=(\Durable,\Rebuild)$ for state whose loss is a downgrade, and
$\mathsf{Retained}=(\Durable,\Retained)$ for custody.  $\mathsf{MemStorage}$ is
$(\Volatile,-)$: admissible only at $\Rebuild$.

\begin{table}[t]
\centering
\caption{Stores of a node and the authority the law assigns}\label{tab:stores}
\footnotesize
\begin{tabular}{@{}lll@{}}
\toprule
Store & Today & Law\\
\midrule
DHT data (native, browser) & $(\Cached,\Rebuild)$ & $(\Cached,\Rebuild)$\\
\quad relay carriers in it & $(\Cached,\Rebuild)$ & $(\Durable,\Retained)$ \issue{1021}\\
Measure, evidence & $(\Cached,\Rebuild)$ & $(\Cached,\Rebuild)$\\
Entry guards & $(\Cached,\Rebuild)$ & $(\Durable,\Rebuild)$ \issue{911}\\
Replay (native) & $(\Durable,\Auth)^{\ast}$ & $(\Durable,\Auth)$ \issue{915}\\
Replay (browser) & $(\Cached,\Rebuild)^{\dagger}$ & $(\Durable,\Auth)$ \issue{912},\issue{923}\\
Replay (no store given) & $(\Volatile,\Auth)$ & refused \issue{1022}\\
\bottomrule
\end{tabular}
\\[2pt]
{\footnotesize $^{\ast}$W missing; D only on unix.  $^{\dagger}$Decodes as $\Auth$ (F) but
neither D (relaxed commit) nor N (row-cap eviction) holds.}
\end{table}

Table~\ref{tab:stores} is the whole classification.  The two rows marked with new issues are
inadmissible points the law exposes outside the five issues: relay carriers are the only
copy of held mail (\issue{1021}), and a provider given no replay store runs replay at
$(\Volatile,\Auth)$ (\issue{1022}).

\section{Presence witness}\label{sec:witness}

\subsection{Choice}

W needs, at restore, a set $\wit\subseteq\Nm$ of names that were durably written, so that
$\wit\setminus\dom\sigma'$ are the lost names.  Fail-closed per stream (\#910) is keyed by
record name: $\mathrm{unavailable}(s)\iff\name(\rec(s))\in\U$.  A witness that does not
name the lost record cannot feed $\U$.  The candidates:

\begin{center}
\footnotesize
\begin{tabular}{@{}lccl@{}}
\toprule
Witness & names loss & $O(1)$/new key & operator step\\
\midrule
count $|\dom\sigma|$ & no & yes & edit count\\
epoch per record & no & yes & edit epoch\\
manifest log & yes & yes & edit binary log\\
per-record marker & yes & yes & remove a file\\
\bottomrule
\end{tabular}
\end{center}

A count or an epoch detects that something is missing but not which stream, so it can only
fail the whole store closed: one lost file would close every stream.  A manifest names the
loss but turns the operator's per-stream recovery into editing a binary log, and its torn
tail needs its own law.  We choose the marker.

\begin{definition}[Presence marker]\label{def:marker}
For an authoritative native store, the marker of name $n$ is an empty file $n.\mathtt{p}$ in
the store's own directory.  Let $R$ be the set of record names listed, $P$ the set of names
with a marker, both from one listing.
\end{definition}

\begin{law}[Write order]\label{law:writeorder}
An authoritative $\mathrm{put}(n,v)$ runs
\begin{align*}
&\mathrm{write}(n.\mathtt{tmp});\ \mathrm{sync};\ \mathrm{rename}(n.\mathtt{tmp},n);\
\mathrm{sync}_{\mathrm{dir}};\\
&[\,n\notin P:\ \mathrm{create}(n.\mathtt{p});\ \mathrm{sync}_{\mathrm{dir}}\,]
\end{align*}
and returns $\Ok$ only after the last step.  An authoritative $\mathrm{rm}(n)$ runs
$\mathrm{remove}(n.\mathtt{p});\mathrm{sync}_{\mathrm{dir}};\mathrm{remove}(n);
\mathrm{sync}_{\mathrm{dir}}$.
\end{law}

\begin{invariant}[Marker implies record]\label{inv:marker}
In every physical state the program and a crash reach, $P\subseteq R$.
\end{invariant}
\begin{proof}
A marker is created only after the rename that makes $n\in R$ is flushed, and removed
before the record.  A crash between any two steps of Law~\ref{law:writeorder} leaves either
$P$ unchanged or the record already durable; by D the flushed prefix is what recovery sees.
\end{proof}

\begin{law}[Restore]\label{law:restore}
Restore classifies, from one listing:
\begin{align*}
n\in R,\ \text{decodes as its key} &\ \mapsto\ \Filed(k,v),\\
n\in R,\ \text{otherwise} &\ \mapsto\ \Undec(n),\\
n\in P\setminus R &\ \mapsto\ \Missing(n),
\end{align*}
and for $n\in R\setminus P$ creates $n.\mathtt{p}$ (flushed) before the store serves its
first operation.  The owner's $\U$ is $\{n:\Undec(n)\}\cup\{n:\Missing(n)\}$.
\end{law}

$R\setminus P$ is the crash between the rename and the marker: the record is whole and
present, so it witnesses itself, and the repair restores $P=R$.

\begin{proposition}[W holds]\label{prop:w}
On a platform that discharges D, Laws~\ref{law:writeorder} and \ref{law:restore} discharge
Obligation~\ref{ob:witness} for every partial loss that leaves markers: a removed or renamed
record file, a record lost by \texttt{fsck}, a rolled back rename of a new record.
\end{proposition}
\begin{proof}
\emph{Complete.}  Let $n\in\dom\sigma_k$; its first $\mathrm{put}$ completed, so
$n\in P$ (Law~\ref{law:writeorder}).  If $n$ is lost and its marker is not, $n\in P\setminus
R$, reported $\Missing$.  \emph{Sound.}  $\Missing(n)$ requires $n\in P$, and a marker
is created only after a put of $n$ flushed its rename (Law~\ref{law:writeorder}), so that
record was durable.  A crash before the marker's creation leaves $n\in R\setminus P$, never
$\Missing$: a put that never reached its rename is never reported lost.
\end{proof}

\begin{proposition}[Total loss is undetectable]\label{prop:total}
No witness stored in the failure domain of the records detects $\lambda_{\mathrm{tot}}$.
\end{proposition}
\begin{proof}
After $\lambda_{\mathrm{tot}}$ the physical state equals that of a store never opened; any
function of the physical state takes the same value on both.
\end{proof}

Total loss is therefore the documented reset: it is what the operator does by removing the
replay store, and its consequences are those of clearing every stream (the recovery law of
\texttt{message::replay}): a receiver may admit an unexpired transaction once more, and a
sender restarts at sequence $0$, which its destinations see as $\mathsf{Stale}$ or
$\mathsf{Fork}$ until it passes their retained high watermark.  A witness outside the domain
(a sibling file in the data directory) would detect a vanished store directory, but the only
producer of that state is the operator, whose clearing step it would then have to name; we
do not add it.

\begin{law}[Read agreement]\label{law:readagree}
An authoritative $\mathrm{get}(k)$ for $\name(k)\in P$ whose record is not readable as the
record of $k$ returns $\Err$, never $\Ok(\bot)$; $\Ok(\bot)$ means $\name(k)\notin P$.
\end{law}

This closes the two cases added to \issue{915} in review: a store whose root vanishes while
the node runs answers $\Err$ for every witnessed name (as its persist already fails), and an
upper-cased record file, which a case-insensitive volume still opens through $\mathrm{get}$,
is $\Missing$ to both $\mathrm{get}$ and $\mathrm{scan}$ because its lower-case marker
remains.

\begin{assumption}[Browsers lose whole origins]\label{as:idbloss}
A browser deletes stored data per origin, never part of it~\cite{MDNEviction}; this
includes storage-pressure eviction of a non-persistent origin and WebKit's deletion of
script-written data after seven days without user interaction.
\end{assumption}

\begin{proposition}[IndexedDB needs no marker]\label{prop:idbwitness}
Under Assumption~\ref{as:idbloss} and N, an IndexedDB store discharges W with $P=\emptyset$.
\end{proposition}
\begin{proof}
Every IndexedDB write is a transaction, atomic by the specification~\cite{IndexedDB3}; with
N the program removes no row it was not told to.  By Assumption~\ref{as:idbloss} the
platform loses no proper subset of a database's rows, so no $\lambda_{\mathrm{part}}$ is
reachable and W holds vacuously; total loss is Proposition~\ref{prop:total}.
\end{proof}



\textbf{Cost.}  W adds one marker creation and one directory flush to the first
$\mathrm{put}$ of each name, and nothing to later puts; restore lists one directory, as
today.  The replay store creates a name once per stream, so the cost is per new stream, not
per transition.

\textbf{Operator step.}  To clear one failed stream, stop the node and remove both
\texttt{<name>} and \texttt{<name>.p} from \texttt{transaction-replay}.  A store written
before W has $P=\emptyset$ and is repaired by Law~\ref{law:restore} at its first open: no
migration is needed.

\section{Survival per platform}\label{sec:survival}

D is a property of the platform's commit, so each platform discharges it separately.

\textbf{Unix.}  Law~\ref{law:writeorder} with
$\mathrm{sync}=\mathtt{File::sync\_all}$, which the Rust standard library maps to
$\mathtt{fcntl}(\mathtt{F\_FULLFSYNC})$ on Apple targets and to $\mathtt{fsync}$
elsewhere~\cite{AppleFsync}, and $\mathrm{sync}_{\mathrm{dir}}$ the same on the directory.
The store's own directory entry is flushed at open, with its ancestors.  The rename-and-sync
pattern is the one crash-consistency studies find portable across unix file
systems~\cite{Pillai2014}.  \texttt{test\_durability} model-checks the step plans over every
crash point.

\textbf{Windows.}  A directory cannot be opened for flushing through the standard library,
so $\mathrm{sync}_{\mathrm{dir}}$ is a no-op today and a completed rename can roll back
(\issue{915}).  The windows gateway is built and run in CI, so the target is in scope.
\begin{law}[Write-through rename]\label{law:writethrough}
On Windows every directory mutation of a $\Durable$ store is a rename by
$\mathtt{MoveFileExW}$ with
$\mathtt{MOVEFILE\_REPLACE\_EXISTING}\mid\mathtt{MOVEFILE\_WRITE\_THROUGH}$, which does not
return until the move is on disk~\cite{MoveFileEx}: a put renames its flushed temporary
file; a marker is created as a flushed temporary file renamed into place; a removal renames
its file to \texttt{.tmp} and then deletes it.
\end{law}
A removal's deletion need not be durable: a \texttt{.tmp} left by a crash is removed at the
next open (the index law of \texttt{FileStorage}), and it is not a record, so it never
reappears.

\textbf{IndexedDB.}  The specification makes the durability hint advisory~\cite{IndexedDB3}.
What each engine does, read from its source:

\begin{center}
\footnotesize\raggedright
\begin{tabular}{@{}p{0.17\columnwidth}p{0.42\columnwidth}p{0.31\columnwidth}@{}}
\toprule
Engine & \jsq{default} & \jsq{strict}\\
\midrule
Chromium & relaxed since 121: complete fires once the data reaches the OS
buffer~\cite{ChromeDurability} & flush before complete\\
Gecko & SQLite WAL, synchronous \texttt{NORMAL} (\texttt{FULL} under a
pref)~\cite{GeckoIDB} & synchronous \texttt{EXTRA}\\
WebKit & SQLite WAL, \texttt{synchronous} left at the library default~\cite{WebKitIDB} &
a \texttt{FULL} WAL checkpoint after commit\\
\bottomrule
\end{tabular}
\end{center}

In WAL mode a commit under \texttt{synchronous=NORMAL} survives a process crash but may roll
back after a system crash~\cite{SQLiteWAL}, which is Chromium's relaxed behaviour too.  So
\jsq{default} discharges D against a process crash in all three engines, and against a
system crash in none that we can rely on.

\begin{law}[Strict commit]\label{law:strict}
Every read-write transaction of an IndexedDB store at $c=\Durable$ is opened with
$\{\mathit{durability}:\text{\jsq{strict}}\}$, and the store reports $\Ok$ only on the
transaction's \texttt{complete} event.
\end{law}

Under Law~\ref{law:strict} D holds to the strength of each engine's strongest commit.  We
claim no more: whether that commit also flushes a drive's write cache is the engine's
choice, as \texttt{F\_FULLFSYNC} is ours on native Apple targets.  \texttt{rexie}~0.4 opens
no transaction with options; the binding is
\texttt{IdbDatabase::\allowbreak transaction\_with\_\allowbreak str\_sequence\_\allowbreak
  and\_mode\_\allowbreak and\_options}, which
\texttt{web-sys} gates behind \texttt{web\_sys\_unstable\_apis}, so the implementation either
enables that cfg for the core crate or passes the options object through \texttt{js-sys}.

\begin{corollary}[Authoritative reads need no write]\label{cor:idbread}
\texttt{IdbStorage} restamps a row on every $\mathrm{get}$ to order eviction.  Under N
eviction never runs, so an authoritative $\mathrm{get}$ may be a read-only transaction.
\end{corollary}

\section{No eviction}\label{sec:noevict}

\begin{law}[Refuse, never evict]\label{law:refuse}
For $m\sqsupseteq\Retained$ on any backend: a $\mathrm{put}$ of a new name into a full store
returns $\Err(\Exhausted)$ with the store unchanged; rewriting an existing name never needs
room; an open with a capacity below the stored count fails.
\end{law}

\begin{proposition}[Eviction is a forged absence]\label{prop:evict}
If an authoritative IndexedDB store evicts, W fails.
\end{proposition}
\begin{proof}
An evicted row $n\in\dom\sigma_k$ is gone with no marker (Proposition~\ref{prop:idbwitness}
has $P=\emptyset$), so restore reports nothing for it and its stream restarts from
$\mathsf{First}$: exactly \issue{923}.
\end{proof}

With Law~\ref{law:refuse} the browser replay store matches the native one (\#914): a new
stream past the bound fails closed, and the row capacity
$\mathtt{TRANSACTION\_REPLAY\_STORE\_MAX\_RECORDS}$ becomes a refusal bound rather than a
sizing hope.  Whether a remote peer can drive a browser node to the cap stops mattering to
safety; reaching it costs only new streams.

\section{Entry guards}\label{sec:guards}

\begin{law}[Guards are a durable cache]\label{law:guards}
The entry-guard store is opened at $(\Durable,\Rebuild)$, \texttt{DurableCache}: writes
discharge D (native: Law~\ref{law:writeorder} without markers, Law~\ref{law:writethrough};
browser: Law~\ref{law:strict}); a guard record that is absent or does not decode is retired
and the set is re-selected from scratch, logged at \texttt{warn} with the record's name.
\end{law}

\begin{proposition}[Guards change only by reconciliation]\label{prop:guards}
Across a crash, the guard set the next selection reads is the last one written, unless a
loss event (Definition~\ref{def:loss}) removed it.
\end{proposition}
\begin{proof}
D gives $\sigma'(n)\in\{\sigma_k(n),\sigma_{k+1}(n)\}$ for the one record per network; both
are whole sets that \texttt{reconcile\_guard\_state} itself produced.
\end{proof}

The meaning stays $\Rebuild$, which answers the policy question of \issue{911}: guards are
privacy state whose loss is a downgrade (one re-selection of the first hop), not a safety
failure, so failing onion routing closed on an undecodable record (F) would trade
availability for nothing the guard set protects.  Under D an undecodable guard record is no
longer reachable by a crash, only by a loss event or a schema change, which the operator
cleans before 1.0.  The browser fallback to an in-memory guard store when IndexedDB cannot
open is $(\Volatile,\Rebuild)$, admissible, and already logged.

\section{Relay custody}\label{sec:custody}

\subsection{Model}

A relay element $e$ carries a signed transaction of stream $s$ (origin account, this node,
class) at sequence $q$ with digest $h$; write $e=(s,q,h)$.  The actors:
\begin{itemize}[leftmargin=1.2em]
\item the carrier $\stored(d)$: the elements held for $d$ at $d+1$, retired only by $d$'s
covering remove up to an element's dot;
\item the receiver window $\W_s$, persisted at $(\Durable,\Auth)$, with the pure transition
$\obs(\W_s,q,h)\in\{\mathsf{First},\mathsf{Advance},\mathsf{Late},\mathsf{Replay},
\mathsf{Fork},\mathsf{Stale}\}$;
\item the application inbox, volatile; \emph{taking} a message is where core's custody
ends.
\end{itemize}
Today a drain admits $h$ into $\W_s$ and queues the message; a restart before the take
loses the queue but not $\W_s$, so the redelivery observes $\mathsf{Replay}$ and is
retired (\issue{936}).  The fact that is lost is ``$h$ was admitted for $e$ and not taken''.

\begin{definition}[Custody set]\label{def:custody}
The custody set of stream $s$ is $\K_s\subseteq\mathbb{N}\times\mathrm{Digest}$: the relay
deliveries of $s$ admitted and not yet taken.  The receiver record of $s$ is the pair
$(\W_s,\K_s)$, one record under the existing key, so one write commits both atomically (A
on both platforms).
\end{definition}

Placing $\K_s$ in the receiver record, and not in the application or a separate store, is
forced by two constraints: core never consults the application (the layering law), and a
second record could not be committed atomically with $\W_s$, which would reopen the window
between admission and its custody.

\begin{law}[Custody]\label{law:custody}
A drain pass first restricts custody to the carrier,
$\K_s\leftarrow\K_s\cap\{(q,h):(s,q,h)\in\stored(d)\}$ (written only if it changed), then
offers each element $e=(s,q,h)$ it has not handed over in this process by
the first arm that applies:
\begin{enumerate}[leftmargin=1.5em]
\item $(q,h)\in\K_s$: $\Requeue$;
\item $|\bigcup_s\K_s|=\mathtt{RELAY\_INBOX\_MAX\_LEN}$: $\Defer$;
\item $\obs(\W_s,q,h)=(\W_s',v)$ with $v$ permitting dispatch: commit
$(\W_s',\K_s\cup\{(q,h)\})$, then queue;
\item otherwise: $\Retire$,
\end{enumerate}
where $\Requeue$ queues the message again without touching $\W_s$, and every queueing
reserves room first ($\Defer$ when there is none, as today).  When the application takes a
message of $e$, core commits $(\W_s,\K_s\setminus\{(q,h)\})$, the \emph{release}; the
element is retired from the carrier only after its release committed.  The live path is
unchanged: it observes $\W_s$ alone.
\end{law}

\begin{lemma}[Bounded custody]\label{lem:bound}
$|\bigcup_s\K_s|\le\mathtt{RELAY\_INBOX\_MAX\_LEN}$ ($64$), so a receiver record grows by
at most $64$ entries of $8+32$ bytes over the whole store.
\end{lemma}
\begin{proof}
The $\Defer$ arm refuses a new entry at the bound, and the restriction step keeps only
entries of elements still stored, of which a carrier holds at most $64$.
\end{proof}

\begin{proposition}[At least once across restart]\label{prop:atleastonce}
Let the carrier satisfy $(\Durable,\Retained)$ (\issue{1021}).  In every history with any
number of crashes, an element $e=(s,q,h)$ leaves the carrier only if (i) the application
took a message of $e$, (ii) the live path would drop $h$ too, or (iii) its retention bound
elapsed.
\end{proposition}
\begin{proof}
The drain retires $e$ in two arms.  After a release: the release follows a take, case (i).
In the last arm of Law~\ref{law:custody}: $(q,h)\notin\K_s$ and $\obs$ refuses.  If $\obs$
is $\mathsf{Replay}$, $h$ is in $\W_s$, so it was admitted earlier either live (dispatched
to the application, case (ii)) or by a relay delivery, which entered $\K_s$ in the same
commit and left it only by a release (case (i)) or by the restriction step, which requires
$e\notin\stored(d)$.  $\mathsf{Fork}$ and $\mathsf{Stale}$ are refusals of the live path,
case (ii).  A crash cannot separate the admission from its custody entry, nor a release
from its effect, because each is one record write under D.  Case (iii) is the retention
law of \#913.
\end{proof}

\begin{proposition}[Duplicate window]\label{prop:dup}
The application receives a message of one element twice iff a crash falls after its take
and before its release committed, and then at most once more per such crash.
\end{proposition}
\begin{proof}
Before the release, $(q,h)\in\K_s$ survives the crash and the redelivery is a
$\Requeue$.  After it, $(q,h)\notin\K_s$ and $h\in\W_s$ or below its window, so the
redelivery is refused and retired.
\end{proof}

The duplicate carries the same transaction id, so an application that needs exactly-once
delivery deduplicates by it.  Moving the release before the take returns would make the
inbox's synchronous \texttt{try\_next} perform a store write; we keep the take synchronous
and accept the window, which is one record write long when the release is issued on the
take's consumption signal.

\subsection{Mail behind live traffic}

\issue{936} also asks that a relay message behind live traffic of the same origin be
delivered.  Inside the window it is: an empty slot gives $\mathsf{Late}$.  Below it, no
receiver can.

\begin{proposition}[No bounded receiver admits stale mail safely]\label{prop:stale}
Let $q<\rmin(\W_s)$ and $(q,h)\notin\K_s$.  No decision that depends only on bounded
receiver state both admits every such $e$ whose $h$ was never dispatched and refuses every
such $e$ whose $h$ was.
\end{proposition}
\begin{proof}
Take two histories that end in the same $(\W_s,\K_s)$: in the first, $h$ arrived live,
was dispatched, and at least $\mathtt{REORDER\_WINDOW}$ later sequences followed; in the
second, $h$ was held, never delivered, and the same later sequences arrived live.  A holder
may hold a payload it forwarded (any hold inside the sender's proof lifetime is
admissible), so $e$ can be stored in both.  The window retains only the last $32$ digests,
so the states agree, and any function of them decides both alike: admitting is a replay in
the first, refusing is a loss in the second.  Separating them needs every digest admitted
within the sender's proof lifetime, which no bound on the store's size covers.
\end{proof}

So below the window the drain refuses as the live path does, which is the criterion
\issue{936} itself states (a refusal retires an element only when core would drop the same
payload arriving live).  Delivering such mail needs the sender to put it on a stream that
live traffic does not advance, a protocol change outside this law.

\section{Implementation order}\label{sec:plan}

\begin{center}
\footnotesize
\begin{tabular}{@{}llp{0.6\columnwidth}@{}}
\toprule
PR & Issue & Law\\
\midrule
1 & \issue{923} & Law~\ref{law:refuse} on \texttt{IdbStorage}; Cor.~\ref{cor:idbread}\\
2 & \issue{912} & Law~\ref{law:strict}; SECURITY.md next to the native text\\
3 & \issue{915} & $\Missing$ in \texttt{ScannedRecord}; Laws~\ref{law:writeorder},
\ref{law:restore}, \ref{law:readagree}, \ref{law:writethrough}\\
4 & \issue{911} & \texttt{DurableCache}; Law~\ref{law:guards}\\
5 & \issue{936} & Law~\ref{law:custody}; receiver record gains $\K_s$\\
\bottomrule
\end{tabular}
\end{center}

The dependencies: 2 and 4 on the browser need the strict-commit seam of 2; 4 on Windows
needs the write-through rename of 3; 5 needs nothing from 1--4, but its
Proposition~\ref{prop:atleastonce} holds in full only with \issue{1021}.  PR~1 comes first because
it is the only one that closes a safety gap reachable without a crash.  \issue{1021} and
\issue{1022} follow the batch.

\textbf{Operator steps}, in the CHANGELOG under \texttt{Unreleased}: PR~5 changes the
receiver record, so an old record does not decode and fails its stream closed; the operator
removes the \texttt{transaction-replay} store (native directory, browser database
\texttt{<storage name>/transaction-replay}) while the node is stopped, accepting the reset of
Proposition~\ref{prop:total}.  PRs 1--4 need no step: an existing native store gains its
markers at its first open, and the IndexedDB schema is unchanged.

\section{Summary of laws}

\begin{center}
\small
\begin{tabular}{@{}ll@{}}
\toprule
Predicate & Discharged by\\
\midrule
presence, native & markers (Laws~\ref{law:writeorder}--\ref{law:readagree})\\
presence, browser & N and whole-origin loss (Prop.~\ref{prop:idbwitness})\\
presence, total loss & undetectable; the documented reset\\
survival, unix & fsync, rename, directory fsync\\
survival, Windows & write-through rename (Law~\ref{law:writethrough})\\
survival, browser & strict commit (Law~\ref{law:strict})\\
presence of custody & $\K_s$ in the receiver record (Law~\ref{law:custody})\\
\bottomrule
\end{tabular}
\end{center}

\sloppy
\bibliographystyle{IEEEtran}
\bibliography{durable-records}

\end{document}
